\heading{Additional practice}{Optional · Not for submission}
The following four exercises are for self-study. They carry \textbf{no points}, are \textbf{not required for submission} and do not affect the Part B grade. The assessed Part B consists only of B1--B3 above. Choose the topics you want to practise.

Quantities are nonnegative, prices and income are positive unless stated otherwise. Show your intermediate steps.
\section*{O1. Representations and group indices}
Let $u(x)=x_1+2x_2$ on $\R_+^2$.
\begin{enumerate}[label=(\alph*)]
\item Which of $3u-7$, $u^2$ and $\min\{u,2\}$ represent the same preferences? Explain using the range of $u$; give two bundles witnessing any failed representation.
\item Draw the indifference sets $u=2$ and $u=4$. Relabel them under the first two transformations without changing the curves.
\item Now consider $F(t,c,z)=(1+\sqrt{tc})(1+z)$, with $t,c,z\ge0$. Write $F=U(v(t,c),z)$ with $U$ strictly increasing in $v$. Explain why the ranking of two $(t,c)$ bundles at a common $z$ does not change when that common $z$ changes. Does this determine a fixed spending budget for the group?
\end{enumerate}
\section*{O2. Budget geometry}
Prices are $p=(2,3)$ and income is $m=18$.
\begin{enumerate}[label=(\alph*)]
\item Write the budget set and budget line. Find both intercepts and the slope, with $x_1$ on the horizontal axis.
\item Check whether $(3,4)$, $(6,3)$ and $(0,5)$ are affordable, and identify which affordable bundles exhaust the budget.
\item In separate diagrams, show what happens when income rises to 36 at unchanged prices, and when $p_1$ rises to 4 at unchanged income 18 and $p_2=3$.
\item Compare the original budget with prices $(4,6)$ and income 36. Then explain why LNS implies budget exhaustion at an existing optimum.
\end{enumerate}

\clearpage
\section*{O3. A kink without derivatives}
Let $u(x)=\min\{x_1,2x_2\}$, with $p_1,p_2,m>0$.
\begin{enumerate}[label=(\alph*)]
\item For a target utility $k\ge0$, write the two inequalities required to attain at least $k$. Find the cheapest bundle attaining it and its cost.
\item Use this cost to find the highest affordable $k$ and derive both Marshallian demands. Explain why the chosen bundle is unique.
\item Find demand and indirect utility at $p=(1,2)$, $m=12$.
\item Why is an ordinary differentiable tangency condition inappropriate at the optimum? Explain how the fixed-proportion argument replaces it.
\end{enumerate}
\section*{O4. Quasilinearity and the income boundary}
Let $u(a,q)=a+5\log(1+q)$, with $a,q\ge0$, both prices equal to 1, and budget $a+q\le m$, $m>0$.
\begin{enumerate}[label=(\alph*)]
\item Explain why utility is quasilinear in $a$ and why the budget is exhausted. Substitute $a=m-q$ and state the feasible interval for $q$.
\item Differentiate the reduced objective and find its unconstrained maximiser. Impose $0\le q\le m$ to obtain demand at all incomes.
\item Calculate the bundle at $m=2$ and $m=6$. Draw a diagram of $q^*(m)$ and explain the change of regime.
\item Derive the indirect utility as a piecewise function of $m$. Check that its value and its first derivative match at the transition.
\item Find the budget multiplier in each regime. Explain when extra income buys more $q$ and when it buys only more $a$.
\end{enumerate}
